\subsection{Modules over Dedekind domains}

We now assume that A is a Dedekind domain; then we know that the ideals \(p \in P\) are maximal and that they are the only prime ideals \(\neq 0\) of A ( \(\S 2\) , no. 1); the group D(A) is identified with the group I(A) of fractional ideals \(\neq 0\) of A.

PROPOSITION 21. Let A be a Dedekind domain. Every pseudo-zero A-module is zero. Every pseudo-injective (resp. pseudo-surjective, pseudo-bijective, pseudo-zero) A-module homomorphism is injective (resp. surjective, bijective, zero).

The first assertion has already been shown (no. 4, Example 1); the others follow from it immediately.

PROPOSITION 22. Let A be a Dedekind domain and M a finitely generated A-module. The following properties are equivalent:

\begin{enumerate}
\def\labelenumi{(\alph{enumi})}
\item
  M is torsion-free.
\item
  M is reflexive.
\item
  M is projective.
\end{enumerate}

We already know (with no hypothesis on the integral domain A) that (b) implies (a) (no. 2, Remark 1) and that (c) implies (b) (Algebra, Chapter II, § 2, no. 7, Corollary 4 to Proposition 14). If M is torsion-free, it is identified with a lattice of \(\mathbf{V} = \mathbf{M} \otimes_{\mathbf{A}} \mathbf{K}\) with respect to A; \(\mathbf{M}_{\mathfrak{p}}\) is therefore a free \(\mathbf{A}_{\mathfrak{p}}\) -module for every maximal ideal \(\mathfrak{p} \in \mathbb{P}\) , since \(\mathbf{A}_{\mathfrak{p}}\) is a principal ideal domain. The conclusion then follows from Chapter II, § 5, no. 2, Theorem 1 (b).

COROLLARY. Let M be a finitely generated A-module and let T be its torsion submodule. Then T is a direct factor of M.

As M/T is torsion-free and finitely generated, it is projective by Proposition 22 and the corollary therefore follows from Algebra, Chapter II, § 2, no. 2, Proposition 4.

PROPOSITION 23. Let A be a Dedekind domain and T a finitely generated torsion A-module. There exist two finite families \((n_{i})_{i\in I}\) and \((\mathfrak{p}_i)_{i\in I}\) , where the \(n_i\) are integers \(\geqslant 1\) and the \(\mathfrak{p}_i\) elements of P, such that T is isomorphic to the direct sum \(\bigoplus_{i\in I} (\mathrm{A} / \mathfrak{p}_{i}^{n_i})\) . Further, the families \((n_{i})_{i\in I}\) and \((\mathfrak{p}_i)_{i\in I}\) are unique to within a bijection of the indexing set.

This follows from no. 4, Theorem 5, taking account of the fact that a pseudo-isomorphism is here an isomorphism.

PROPOSITION 24. Let A be a Dedekind domain and M a finitely generated torsion-free A-module of rank \(n \geqslant 1\) . Then there exists an ideal \(d \neq 0\) of A such that M is isomorphic to the direct sum of the modules \(A^{n-1}\) and d.~Moreover, the class of the ideal d is determined uniquely by this condition.

Theorem 6 of no. 9 shows that there exists a free submodule L of M such that M/L is isomorphic to an ideal a of A. If a = 0, we take d = A. Otherwise, a is of rank 1, hence \(L = A^{n-1}\) and a is a projective module (Proposition 22); M is therefore isomorphic to the direct sum of L and a (Algebra, Chapter II, §2, no. 2, Proposition 4), which proves the first part of the proposition. Moreover, it follows from no. 7, Proposition 16 (i), (iv) and (v) that \(c(\mathbf{M}) = c(\mathfrak{d})\) whence the uniqueness of the class of d.

\textbf{Remarks}\par

\begin{enumerate}
\def\labelenumi{(\arabic{enumi})}
\item
  Propositions 23 and 24 and the Corollary to Proposition 22 determine completely the structure of finitely generated A-modules. Proposition 24 shows that a torsion-free finitely generated A-module is determined up to isomorphism by its rank and the divisor class which is attached to it.
\item
  It can be shown that over a Dedekind domain a projective module which is not finitely generated is necessarily free (Exercise 21) and that every sub-module of a projective module is projective (Exercise 20).
\end{enumerate}


